\documentclass[./dissertation.tex]{subfiles}
\begin{document}

\chapter{Related Work}

\section{Learning Success Models for Grasping}
Robotic grasping is a major area of study within robotics due to its wide use in manipulation tasks. As these tasks become more complex, traditional hand-engineered approaches are becoming increasingly ineffective. As a result, a substantial portion of research in robotic grasping has shifted toward learning-based methods that leverage datasets to model grasp outcomes. Many of these approaches focus on learning a grasp success model that predicts whether a proposed grasp action will succeed in the current state of the environment.

These models are trained to predict the success probability of candidate grasps using information such as RGB images, depth maps, object positions, or affordance representations. For example, after collecting data from 50,000 grasp attempts, \parencite{SSS} trained a Convolutiona Neural Network (CNN) that takes image patches as input and outputs the probability of grip success for a set of candidate grasp angles. Similarly, \parencite{LevineGrasp} learned a grasp prediction network that predicts the likelihood of a successful grasp from visual input and integrated this estimate within a continuous visual servoing controller. 

Other work has explored the use of synthetic datasets to train these grasp evaluation models. After generating a large dataset of synthetic point clouds, \parencite{DexNet} learned a grip robustness function that evaluates candidate grasps and enables planning by selecting the highest-scoring action. This approach demonstrates how simulated data can be used to train these models while allowing them to generalize to real world tasks. Other methods, such as \parencite{DSGD} and \parencite{detry2011learning}, learn these grasp success models from interaction data and utilize them for simple evaluation or to find high success actions.

The method in this paper follows this line of work by learning a grasp success model that maps states and actions to success probabilities. However, rather than using this learned model solely as an evaluation metric for grasp selection, the success predictor is interpreted as a reward function in this work. This reward model is then incorporated into a reinforcement learning framework to optimize an effective grasping policy.

\section{Robust Reinforcement Learning}
Robust Reinforcement Learning \parencite{RobustRLSurvey} encompasses a class of methods that extend the standard reinforcement learning objective to account for hostile environmental effects or adversarial factors. These approaches are especially useful for real world applications, where uncertainty and disturbances are inherent to the environment. The goal therefore is to learn a policy that achieves high expected reward while remaining effective in the presence of noise, adversaries, and uncertainty.

Prior work in Robust RL can be categorized based on the source of uncertainty that the method aims to address. The most common approaches focus on robustness with respect to uncertainty in the transition dynamics or reward function. For example, \parencite{RobustTransition} accounts for the stochastic nature of the environment by reformulating the reinforcement learning objective as a minimax problem, where the policy is optimized against the worst-case dynamics drawn from the set of possible transition probabilities. Additionally, \parencite{RobustReward} learns a policy that maximizes a combination of a performance and robustness measure, where the robustness component captures the worst-case reward for the evaluated action. By optimizing against the worst-case outcomes in the reward or transition dynamics, these approaches produce policies that are robust but often conservative, sometimes resulting in reduced performance.

Many formulations that address transition and reward uncertainty are limited to tabular settings and do not easily extend to nonlinear function approximators or continuous state-action spaces. Another line of research in Robust RL instead focuses on robustness to adversarial or noisy actions, which can be applied in these more complex settings. The work of \parencite{RobustActionMDP} introduces the Noisy Action Robust MDP, where the standard policy is replaced with a joint policy $\pi_{mix}$ that perturbs the intended action through an adversarial policy. The resulting optimization is formulated in a minimax framework for the mixed policy. Similarly, \parencite{RobustActionDelta} defines a robust optimization objective in which adversarial perturbations are applied to actions in order to minimize the expected reward before policy optimization occurs. Much like the transition and reward methods, these methods aim to learn policies that perform well under worst-case action perturbations.

The objective proposed in this work draws a significant amount of inspiration from Robust RL techniques. However, rather than optimizing for the worst-case scenario like the methods described above, our approach instead optimizes with respect to a fixed perturbation magnitude to the actions. This formulation aims to produce policies that remain effective under environmental uncertainty without becoming overly conservative, thereby balancing robustness and performance.

\section{Offline Reinforcement Learning}
Offline Reinforcement Learning \parencite{OfflineRLSurvey} focuses on learning a policy $\pi$ entirely from a fixed dataset $\mathcal{D}$ that has been previously collected by an exploration or behavior policy $\pi_\beta$, without any additional interaction with the environment (see Figure~\ref{Off RL}). This approach contrasts with Online Reinforcement Learning, where the agent continually collects new experience while learning. Depite this distinction, the main objective of Offline RL is the same as standard Online RL, which is to learn a policy that maximizes expected cumulative reward. Because training is performed on static, batched data, Offline RL is often more sample-efficient and can converge faster than Online RL, where interactions with the environment may be costly and slow down the learning process. 

However, a central challenge with Offline RL is distribution shift. During training, the learned policy $\pi$ may select actions or reach states that were rarely or never encountered in the dataset $\mathcal{D}$. This mismatch between the training data distribution and the states and actions encountered by the learned policy can result in unreliable value estimation, which in turn may lead to poor policy performance when deployed. To address this issue, prior work has developed a range of strategies.Techniques such as importance sampling \parencite{IS} help correct the bias induced by off-policy evaluation, while regularization methods \parencite{Offline2} can encourage desirable behaviors in the learned policy that may not be sufficiently represented in the dataset $\mathcal{D}$. Similarly, policy constraints \parencite{Offline1} can be used to keep the learned policy $\pi$ close to the dataset $\mathcal{D}$, improving the reliability of value estimation. 

Although the method presented in this paper incorporates principles from Offline RL, it does not implement a pure offline approach. The learned policy $\pi$ is subsequently deployed in the environment to acquire more data, which is then used to iteratively refine the reward model for further offline training. Despite this, the strategy of regularization formed a primary inspiration for this work. The combination of offline training with online data collection gives this method the efficiency of Offline RL while also addressing the limitations of static datasets.

\begin{figure}[h]
\centering\includegraphics[width=0.75\textwidth]{figures/Off RL.png}
\caption{Offline Reinforcement Learning. An exploration policy $\pi_\beta$ interacts with the environment to collect a dataset of state-action transitions and rewards, which is stored in the buffer $\mathcal{D}$. This dataset is used to train a policy $\pi$, which is then deployed in the environment without further interaction.}
\label{Off RL}
\end{figure}

\section{Model-Based Reinforcement Learning}
Model-Based Reinforcement Learning \parencite{MBRLSurvey} is a class of methods that learn a model of the environment's dynamics or transition function, which predicts the resulting state of actions taken by the agent. The learned transition model allows the agent to simulate future states and generate synthetic data without interacting with the real environment. This is particularly useful in domains where real world interactions are costly, dangerous, or time consuming. Leveraging a learned model, these methods significantly improve sample efficiency since the agent can perform planning and optimization on a combination of real and simulated experiences. A basic Model-Based RL pipeline is illustrated in Figure~\ref{MBRL}.

A key challenge in Model-Based RL is model accuracy and uncertainty. Several approaches address this issue through uncertainty-aware modeling, such as ensemble methods which train multiple independent models. One approach, proposed in \parencite{MBRL1}, trains an ensemble of $K$ dynamics models and randomly selects one for each rollout, which helps reduce overfitting and increases robustness to individual errors. Additionally, \parencite{MBRL2} learns $N$ reward models and estimates the reward as the ensemble mean penalized by the maximum standard deviation to incentivize the policy to avoid regions of high uncertainty. 

In addition to ensemble methods, many Model-Based RL methods adopt iterative model refinement and Dyna-style \parencite{Dyna} techniques to continuously improve the dynamics model and policy. Works such as \parencite{MBRL3} and \parencite{MBRL4} follow this paradigm by alternating between model learning and data collection with the optimized policy to progressively improve performance. This leads to more robust and effective agents, especially in complex environments. 

Although the proposed approach is related to Model-Based RL, it does not explicitly model a transition function. Instead, the method focuses on learning a reward model, which estimates the expected return of actions directly. Using the reward predictions, the policy can be optimized without the need for modeling state transitions. Despite this distinction, the approach adopts ideas from Model-Based RL, including the use of ensemble networks to capture uncertainty and iterative model improvement to enhance model accuracy.

\begin{figure}[h]
\centering\includegraphics[width=0.85\textwidth]{figures/MBRL.png}
\caption{Model-Based Reinforcement Learning. A policy $\pi_k$ is executed in the real world to collect data, which populates the real world buffer $\mathcal{D}$. This data is used to learn an approximate transition function $\hat{f}_k$. The learned transition function is used to generate a simulated dataset $\hat{\mathcal{D}}$, which is combined with the real-world dataset $\mathcal{D}$ to update the policy $\pi_{k+1}$. This iterative process improves sample efficiency by utilizing simulated data while grounding learning in real world experience.}
\label{MBRL}
\end{figure}

\end{document}
